\subsection{Coxeter 变换的特征值}

由于所有 Coxeter 变换彼此共轭（no. 1, Prop. 1），它们都有同一特征多项式 P(T)。设 h 是 W 的 Coxeter 数。则

\[
\mathrm{P} (\mathrm{T}) = \prod_ {j = 1} ^ {l} \left(\mathrm{T} - \exp \frac {2 i \pi m _ {j}}{h}\right),
\]

其中 \(m_{1}, m_{2}, \ldots, m_{l}\) 是满足下式的整数

\[
0 \leqslant m _ {1} \leqslant m _ {2} \leqslant \dots \leqslant m _ {l} <   h.
\]

定义 2. 整数 \(m_1, m_2, \ldots, m_l\) 称为 W 的指数。

设 C 是由 \(\mathfrak{H}\) 确定的一个室，\(\mathrm{H}_1, \ldots, \mathrm{H}_l\) 是它的诸墙，并置 \(s_i = s_{\mathrm{H}_i}\)。用 \(e_i\) 表示与 \(\mathrm{H}_i\) 正交且与 C 位于 \(\mathrm{H}_i\) 同一侧的单位向量。由 Chap. IV 附录的 Prop. 2，可以假设诸 \(\mathrm{H}_i\) 如此编号，使得 \(e_1, e_2, \ldots, e_r\) 两两正交，且 \(e_{r+1}, e_{r+2}, \ldots, e_l\) 两两正交。于是 \(s' = s_1 s_2 \ldots s_r\) 是关于子空间

\[
\mathrm{V} ^ {\prime} = \mathrm{H} _ {1} \cap \mathrm{H} _ {2} \cap \dots \cap \mathrm{H} _ {r},
\]

\(s^{\prime \prime} = s_{r + 1}s_{r + 2}\dots s_{l}\) 是关于子空间

\[
\mathrm{V} ^ {\prime \prime} = \mathrm{H} _ {r + 1} \cap \mathrm{H} _ {r + 2} \cap \dots \cap \mathrm{H} _ {l},
\]

的正交对称，而 \(c = s's''\) 是一个 Coxeter 变换。由于 \((e_{1}, \ldots, e_{l})\) 是 V 的基，V 是 \(V'\) 与 \(V''\) 的直和。

我们首先推出 1 不是 c 的特征值。因为若 \(x \in V\) 满足 \(c(x) = x\)，则 \(s'(x) = s''(x)\)，于是 \(x - s'(x) = x - s''(x)\) 与 \(V'\) 和 \(V''\) 都正交，从而为零。因此 \(x = s'(x) = s''(x) \in V' \cap V'' = \{0\}\)。

从而

\[
0 <   m _ {1} \leqslant m _ {2} \leqslant \dots \leqslant m _ {l} <   h.\tag{1}
\]

c 的特征多项式具有实系数。因此，对所有 j，\(T - \exp\frac{2i\pi m_{j}}{h}\) 在 \(P(T)\) 中的幂次等于 \(T - \exp\frac{2i\pi(h - m_{j})}{h}\) 的幂次。由此得

\[
m _ {j} + m _ {l + 1 - j} = h \quad (1 \leqslant j \leqslant l).\tag{2}
\]

把等式 (2) 相加，我们得到

\[
m _ {1} + m _ {2} + \dots + m _ {l} = \frac {1}{2} l h.\tag{3}
\]

引理 2. 假设 W 不可约且 \(\dim V \geqslant 2\)。沿用前面的记号，存在两个线性无关的向量 \(z', z''\)，使得

\begin{enumerate}
\def\labelenumi{(\roman{enumi})}
\item
  由 \(z', z''\) 生成的平面 P 在 \(s'\) 和 \(s''\) 下稳定；
\item
  \(s'|\mathrm{P}\) 与 \(s''|\mathrm{P}\) 是关于 \(\mathbf{R}z'\) 与 \(\mathbf{R}z''\) 的正交反射。
\item
  \(z', z'' \in \overline{\mathbf{C}}\)，且 \(P \cap C\) 是 \(z', z''\) 的系数 \(> 0\) 的线性组合所成的集合。
\end{enumerate}

设 \((e^{1},\ldots,e^{l})\) 是 V 中满足 \((e^{i}|e_{j})=\delta_{ij}\) 的基。于是 C 是由诸 \(e^{i}\) 确定的开单纯锥（\(\S 3\)，no. 9, Prop. 7）。显然 \(V'\) 由 \(e^{r+1},\ldots,e^{l}\) 生成，而 \(V''\) 由 \(e^{1},\ldots,e^{r}\) 生成。设 q 是 V 的自同态，满足 \(q(e^{1})=e_{1},\ldots,q(e^{l})=e_{l}\)。它关于 \((e^{1},\ldots,e^{l})\) 的矩阵是 \(Q=((e_{i}|e_{j}))\)。我们有 \((e_{i}|e_{j})\leqslant0\)（对 \(i\neq j\)）（\(\S 3\)，no. 4, Prop. 3）。由于 W 不可约，不存在任何分划 \(\{1,2,\ldots,l\}=I_{1}\cup I_{2}\) 使得 \((e_{i}|e_{j})=0\) 对 \(i\in I_{1}\) 且 \(j\in I_{2}\) 成立。于是（\(\S 3\)，no. 5, Lemma 4）Q 有一个特征向量 \((a_{1},\ldots,a_{l})\)，其所有坐标都 \textgreater0；设 a 是相应的特征值。置

\[
\begin{array}{r} z = a _ {1} e ^ {1} + \dots + a _ {l} e ^ {l}, \\ z ^ {\prime \prime} = a _ {1} e ^ {1} + \dots + a _ {r} e ^ {r} \in V ^ {\prime \prime} \cap \overline {{C}}, \\ z ^ {\prime} = a _ {r + 1} e ^ {r + 1} + \dots + a _ {l} e ^ {l} \in V ^ {\prime} \cap \overline {{C}}, \end{array}
\]

并设 P 是由 \(z'\) 与 \(z''\) 生成的平面。则 \(P \cap C\) 是 \(z'\) 与 \(z''\) 的系数 \(> 0\) 的线性组合所成的集合。关系式 \(q(z) = az\) 给出 \(\sum_{j=1}^{l} a_j e_j = \sum_{j=1}^{l} aa_j e^j\)；用 \(e_k\)（其中 \(k \leqslant r\)）作数量乘法得 \(a_k + \sum_{j=r+1}^{l} a_j (e_j | e_k) = aa_k\)；于是

\[
\begin{array}{l} (a - 1) z ^ {\prime \prime} = \sum_ {k = 1} ^ {r} \left(\sum_ {j = r + 1} ^ {l} a _ {j} (e _ {j} | e _ {k})\right) e ^ {k} \\ \qquad = \sum_ {j = r + 1} ^ {l} a _ {j} \left(\sum_ {k = 1} ^ {r} (e _ {j} | e _ {k}) e ^ {k}\right) \\ \qquad = \sum_ {j = r + 1} ^ {l} a _ {j} \left(- e ^ {j} + \sum_ {k = 1} ^ {l} (e _ {j} | e _ {k}) e ^ {k}\right) \\ \qquad = - \sum_ {j = r + 1} ^ {l} a _ {j} e ^ {j} + \sum_ {j = r + 1} ^ {l} a _ {j} e _ {j} \\ \qquad = - z ^ {\prime} + \sum_ {j = r + 1} ^ {l} a _ {j} e _ {j}. \end{array}
\]

于是，\((a-1)z'' + z'\) 与 \(e^{1}, \ldots, e^{r}\) 正交，即与 \(V''\) 正交。因此，\(s''\) 使由 \(z''\) 与 \((a-1)z'' + z'\) 生成的平面稳定，该平面即 P。同理，\(s'\) 使 P 稳定。由于 \(z' \in P \cap V'\) 且 \(z'' \in P \cap V''\)，\(s'|P\) 与 \(s''|P\) 是关于 \(Rz'\) 与 \(Rz''\) 的反射。

定理 1. 假设 W 不可约。则：

\begin{enumerate}
\def\labelenumi{(\roman{enumi})}
\item
  \(m_{1} = 1, m_{l} = h - 1\) .
\item
  \(\operatorname{Card}(\mathfrak{H}) = \frac{1}{2} lh\) .
\end{enumerate}

我们沿用前面的记号。\(c = s's''\) 在 \(P\) 上的限制是角为 \(2(z'', z')\) 的旋转（§2, no. 5, Cor. of Prop. 6）。由于 \(c\) 的阶为 \(h\)，\(h\) 个元素 \(1, c, \ldots, c^{h-1}\)（均属于 \(W\)）两两不同；于是诸元素 \(s', s'c, \ldots, s'c^{h-1}\) 也两两不同，并且不同于前面的诸元素，因为 \(c^i |P\) 是旋转而 \(s'c^j |P\) 是反射。集合

\[
\{1, c, \dots , c ^ {h - 1}, s ^ {\prime}, s ^ {\prime} c, \dots , s ^ {\prime} c ^ {h - 1} \}
\]

是 W 的子群 \(W'\)，它由 \(s'\) 与 \(s''\) 生成，并且在 P 上诱导出群 \(W''\)，后者由关于 \(Rz', Rz''\) 的正交反射生成。C 在 \(W'\) 的元素之下的像，或与 -C 不相交，或等于 -C。于是，\(P \cap C\) 在 \(W''\) 的元素之下的像，或与 \(-(P \cap C)\) 不相交，或等于 \(-(P \cap C)\)。因此，对 \(P\) 的适当定向，存在整数 \(m > 0\) 使得 \((z'', z') = \frac{\pi}{m}\)（§2, no. 5, Cor. of Prop. 7）。此外，诸集合 \(g'(C)\)（对 \(g' \in W'\)）两两不相交；于是诸集合 \(g''(P \cap C)\)（对 \(g'' \in W''\)）也两两不相交；故 \(W''\) 的阶为 \(2h\)。从而 \(m = h\)。按定义，\(c|P\) 是角为 \(\frac{2\pi}{m}\) 的旋转，故以 \(\exp \frac{2i\pi}{h}, \exp \frac{2i\pi(h - 1)}{h}\) 为特征值。这就证明了 \(m_1 = 1, m_l = h - 1\)。

\(\mathbf{R}z'\) 与 \(\mathbf{R}z''\) 在 \(W'\) 之下的像是 \(h\) 条直线 \(D_1, \ldots, D_h\)（在 \(P\) 中），而 \(P - (D_1 \cup \cdots \cup D_h)\) 的点是 \(W'\) 的元素作用于 \(P \cap C\) 的点所得的像。于是，\(\mathfrak{H}\) 中的一个超平面必然切割 \(P\) 且沿直线 \(D_i\) 之一，从而是一个超平面在 \(W'\) 的一个操作之下的像，该超平面属于 \(\mathfrak{H}\) 且包含 \(\mathbf{R}z'\) 或 \(\mathbf{R}z''\)。

现在，任何 \(H \in \mathfrak{H}\) 只要包含 \(Rz'\)，就都是超平面 \(H_{1}, \ldots, H_{r}\) 之一。事实上，设 \(e_{H}\) 是与 H 正交且与 C 位于 H 同一侧的单位向量。则 \(e_{H} = \lambda_{1}e_{1} + \cdots + \lambda_{l}e_{l}\)，其中诸 \(\lambda_{i}\) 都 \(\geqslant 0\)（§3, no. 5, Lemma 6, (i)）。而 \(0 = (e_{\mathrm{H}}|z') = \lambda_{r+1}a_{r+1} + \cdots + \lambda_{l}a_{l}\)，故

\[
\lambda_ {r + 1} = \dots = \lambda_ {l} = 0, \quad \text { and } \quad e _ {\mathrm{H}} = \lambda_ {1} e _ {1} + \dots + \lambda_ {r} e _ {r}.
\]

假设诸 \(\lambda_{i}\) 中有两个非零，例如 \(\lambda_{1}\) 与 \(\lambda_{2}\)；由于 \(e_{1},\ldots,e_{r}\) 两两正交，我们将有

\[
s _ {1} (e _ {\mathrm{H}}) = - \lambda_ {1} e _ {1} + \lambda_ {2} e _ {2} + \dots + \lambda_ {r} e _ {r}
\]

于是 \(s_{1}(e_{\mathrm{H}})\) 的坐标将不同号，这是荒谬的（loc. cit.）。因此，\(e_{H}\) 与向量 \(e_{1},\ldots,e_{r}\) 之一成比例，这就证明了我们的断言。同理，任何 \(H\in \mathfrak{H}\) 只要包含 \(Rz''\)，就都是超平面 \(H_{r+1},\ldots,H_{l}\) 之一。

于是，\(\mathfrak{H}\) 中包含 \(\mathbf{R}z'\) 或 \(\mathbf{R}z''\) 的元素个数为 \(l\)。若 \(h\) 为偶数，则 \(\operatorname{Card}(\mathfrak{H})\) 等于 \(\frac{h}{2} l\)。若 \(h\) 为奇数，则 \(\operatorname{Card}(\mathfrak{H})\) 等于 \(\frac{h - 1}{2} l + r\)，同时也等于 \(\frac{h - 1}{2} l + (l - r)\)；故 \(r = l - r\)，即 \(r = \frac{l}{2}\)，从而 \(\operatorname{Card}(\mathfrak{H}) = \frac{h - 1}{2} l + \frac{l}{2} = \frac{h}{2} l\)。

注。沿用前面证明中的记号。设 \(c'\) 是 \(\mathbf{C}\)-线性扩张，即 \(c\) 到 \(V \otimes_{\mathbf{R}} \mathbf{C}\) 的扩张，\(c''\) 是 \(c'\) 在 \(P \otimes_{\mathbf{R}} \mathbf{C}\) 上的限制。由我们对 \(c|P\) 的研究可知，\(c''\) 有一个特征向量 \(x\)，它对应于特征值 \(\exp \frac{2i\pi}{h}\)，且该特征向量不属于任何集合 \(D \otimes_{\mathbf{R}} \mathbf{C}\)，其中 \(D\) 表示 \(P\) 的一条直线（因为 \(D\) 在 \(c\) 下不稳定）。而由前文，对任何 \(H \in \mathfrak{H}\)，\(H \cap P\) 是一条直线；故 \(x \notin H \otimes_{\mathbf{R}} \mathbf{C}\)。

系。设 \(R_{0}\) 是 V 中与 H 的某个元素正交的单位向量所成的集合。若 W 不可约，则

\[
\sum_ {u \in R _ {0}} (x | u) ^ {2} = h (x | x)\tag{4}
\]

对所有 \(x \in V\) 成立。

置 \(f(x) = \sum_{u \in R_0} (x|u)^2\)。显然 \(f\) 是一个在 \(W\) 下不变的正二次型，且由于诸 \(e_i\) 构成 \(V\) 的基，它是非退化的。由于

W 不可约，故存在常数 \(\beta\) 使得 \(f(x) = \beta (x|x)\)（§ 2, no. 1, Prop. 1）。若 \((x_{i})_{1\leqslant i\leqslant l}\) 是 V 关于内积 \((x|y)\) 的标准正交基，则

\[
\begin{array}{l} \beta l = \sum_ {i = 1} ^ {l} \beta (x _ {i} | x _ {i}) = \sum_ {i = 1} ^ {l} f (x _ {i}) = \sum_ {i = 1} ^ {l} \sum_ {u \in R _ {0}} (x _ {i} | u) ^ {2} \\ = \sum_ {u \in R _ {0}} 1 = \operatorname{Card} (R _ {0}) = 2 \operatorname{Card} (\mathfrak {H}) = h l. \end{array}
\]

于是 \(\beta = h\)，这就证明了 (4)。

命题 2. 若 \(W\) 不可约且 \(h\) 为偶数，则 \(W\) 中把 \(C\) 变到 \(-C\) 的唯一元素是 \(c^{h/2}\)。

我们使用 Th. 1 证明中的记号。由于 \(c|P\) 是角为 \(\frac{2\pi}{h}\) 的旋转，\(c^{h/2}\) 把 \(z'\) 变到 \(-z'\)，把 \(z''\) 变到 \(-z''\)，从而把 \(z' + z'' = z\) 变到 \(-z\)。而 \(z \in C\)，故室 \(c^{h/2}(C)\) 必为 \(-C\)。

命题 3. 假设 W 不可约。设 \(u_1, \ldots, u_l\) 是对称代数 \(S = \mathbf{S}(V)\) 的齐次元素，它们在 \(\mathbf{R}\) 上代数无关，且生成 \(S\) 中在 \(W\) 下不变的元素所成的代数（§5, no. 3, Th. 3）。若 \(p_j\) 是 \(u_j\) 的次数，则 \(W\) 的指数为 \(p_1 - 1, \ldots, p_l - 1\)。

置 \(V' = V \otimes_{\mathbf{R}} \mathbf{C}\)，\(S' = \mathbf{S}(V') = S \otimes_{\mathbf{R}} \mathbf{C}\)，并把 \(V\) 上的内积扩张为 \(V'\) 上的一个 Hermite 型。若 \(c\) 是 \(W\) 的一个 Coxeter 变换，则存在标准正交基 \((X_i)_{1 \leqslant i \leqslant l}\)，它由 \(V'\) 中 \(c \otimes 1\) 的特征向量构成（Algebra, Chap. IX, §7, no. 3, Prop. 4）；此外，可以假设对 \(1 \leqslant j \leqslant l\)，\(X_j\) 对应于特征值 \(\exp \frac{2i\pi m_j}{h}\)（\(c \otimes 1\) 的特征值）。显然 \(S'\) 可与代数 \(\mathbf{C}[X_1, \ldots, X_l]\) 等同，并且可以写 \(u_j \otimes 1 = f_j(X_1, \ldots, X_l)\)，其中 \(f_j\) 是次数为 \(p_j\) 的齐次多项式（在 \(\mathbf{C}[X_1, \ldots, X_l]\) 中）。置 \(D_j = \frac{\partial}{\partial X_j}\) 与 \(J = \det(D_k f_j)\)。回顾（§5, no. 4, Prop. 5），\(J(X_1, \ldots, X_l)\) 与 \(S'\) 中 \(\text{Card}(\mathfrak{H})\) 个向量 \(y_k\)（属于 \(V\)）的乘积成比例，其中每个向量都与 \(\mathfrak{H}\) 的一个超平面正交。由于可以假设 \(X_1 \notin H \otimes C\) 对所有 \(H \in \mathfrak{H}\) 成立（注），\(X_1\) 在每个向量 \(y_k\) 中的分量均非零，故 \(J(1, 0, \ldots, 0) \neq 0\)。于是由行列式的展开法则可知，存在一个置换 \(\sigma\)（\(\{1, 2, \ldots, l\}\) 的置换），使得 \((D_{\sigma(j)} f_j)(1, 0, \ldots, 0) \neq 0\) 对所有 \(j\) 成立。由于 \(D_{\sigma(j)} f_j\) 是次数为 \(p_j - 1\) 的齐次多项式，\(X_1^{p_j - 1} X_{\sigma(j)}\) 在 \(f_j(X_1, \ldots, X_l)\) 中的系数非零。而 \(f_j(X_1, \ldots, X_l)\) 在 \(c \otimes 1\) 下不变，且

\[
(c \otimes 1) (\mathrm{X} _ {1} ^ {p _ {j} - 1} \mathrm{X} _ {\sigma (j)}) = \left(\exp \frac {2 i \pi}{h} (p _ {j} - 1 + m _ {\sigma (j)})\right) (\mathrm{X} _ {1} ^ {p _ {j} - 1} \mathrm{X} _ {\sigma (j)}).
\]

这就证明 \(p_j - 1 + m_{\sigma(j)} \equiv 0 \pmod{h}\)。而 \(h - m_{\sigma(j)}\) 是一个指数（公式 (2)）。必要时置换诸 \(u_j\)，可以假设 \(p_j - 1 \equiv m_j \pmod{h}\) 对所有 \(j\) 成立。由于 \(p_j - 1 \geqslant 0\) 且 \(m_j < h\)，有 \(p_j - 1 = m_j + \mu_j h\)，其中 \(\mu_j\) 是 \(\geqslant 0\) 的整数。由 §5, Prop. 3 可知

\[
\operatorname{Card} (\mathfrak {H}) = \sum_ {j = 1} ^ {l} (p _ {j} - 1) = \sum_ {j = 1} ^ {l} m _ {j} + h \sum_ {j = 1} ^ {l} \mu_ {j}.
\]

把公式 (3) 与 Th. 1 (ii) 结合起来，得到 \(h \sum_{j=1}^{l} \mu_j = 0\)，故 \(\mu_j = 0\) 对所有 \(j\) 成立，最终得到 \(p_j - 1 = m_j\) 对所有 \(j\) 成立。

系 1. 若 \((m_i)_{1\leqslant i\leqslant l}\) 是 W 的指数的递增序列，则 W 的阶等于 \((m_1 + 1)(m_2 + 1)\ldots (m_l + 1)\)。

这由关系 \(m_j + 1 = p_j\) 以及 §5, Cor. of Th. 3 得出。

系 2. 若 \(c\) 是 \(W\) 的一个 Coxeter 变换，则

\[
\exp \left(\frac {2 i \pi}{h}\right) \quad \text { and } \quad \exp \left(\frac {- 2 i \pi}{h}\right)
\]

都是 \(c\) 的重数为 1 的特征值。

否则，S 中将存在两个不成比例的 2 次齐次不变元，从而存在两个在 W 下不变的 V* 上不成比例的二次型，这与 §2, no. 1, Prop. 1 相矛盾。

系 3. V 的比为 -1 的位似属于 W 当且仅当 W 的所有指数都是奇数。此时 h 为偶数，且对 W 的任何 Coxeter 变换 c 都有 \(c^{h/2} = -1\)。

第一个断言由 §5, no. 3, Prop. 4 得出。假设 \(W\) 的指数都是奇数。则由公式 (2)，\(h\) 为偶数，且

\[
\left(\exp \frac {2 i \pi m _ {j}}{h}\right) ^ {h / 2} = \exp (i \pi m _ {j}) = - 1;
\]

于是 \(c^{h/2} = -1\)，因为 \(c\) 是 \(V\) 的半单自同构（Algebra, Chap. IX, §7, no. 3, Prop. 4）。

